% 可信机器学习研究地图：中心问题、三类条件、八个研究问题与十五项具体课题。
\begingroup
\newcommand{\trustsector}[5]{%
  \path[fill=#5,draw=black,line width=.75pt]
    (#1:#4) arc[start angle=#1,end angle=#2,radius=#4]
    -- (#2:#3) arc[start angle=#2,end angle=#1,radius=#3] -- cycle;}
\newcommand{\trustlabel}[5][]{%
  \node[font=\figurefont\mdseries#1,text=FigureInk,align=center,
    rotate=#5,transform shape] at (#2:#3) {#4};}

\begin{tikzpicture}[figure picture,foundation picture]
  \path[use as bounding box] (-82,-82) rectangle (82,82);

  % 第一层：要求与影响、证据与行为、防护与治理。
  \trustsector{90}{-45}{16}{34}{FigureBlue!78!white}
  \trustsector{-45}{-180}{16}{34}{FigureYellow!82!white}
  \trustsector{-180}{-270}{16}{34}{FigureRed!78!white}

  % 第二层：八个相互制约的研究问题。
  \trustsector{90}{45}{34}{55}{FigureBlue!58!white}
  \trustsector{45}{0}{34}{55}{FigureBlue!58!white}
  \trustsector{0}{-45}{34}{55}{FigureBlue!58!white}
  \trustsector{-45}{-90}{34}{55}{FigureYellow!62!white}
  \trustsector{-90}{-135}{34}{55}{FigureYellow!62!white}
  \trustsector{-135}{-180}{34}{55}{FigureYellow!62!white}
  \trustsector{-180}{-225}{34}{55}{FigureRed!58!white}
  \trustsector{-225}{-270}{34}{55}{FigureRed!58!white}

  % 第三层：每个类别内部顺时针展开，末项在中层收束。
  \foreach \a/\b in {90/75,75/60,60/45,45/30,30/15,15/0}
    {\trustsector{\a}{\b}{55}{80}{FigureBlue!34!white}}
  \foreach \a/\b in {-45/-60,-60/-75,-75/-90,-90/-105,-105/-120,-120/-135}
    {\trustsector{\a}{\b}{55}{80}{FigureYellow!38!white}}
  \foreach \a/\b in {-180/-195,-195/-210,-210/-225}
    {\trustsector{\a}{\b}{55}{80}{FigureRed!34!white}}

  % 中心问题。
  \fill[white] (0,0) circle[radius=16];
  \draw[draw=black,line width=1pt] (0,0) circle[radius=16];
  \node[font=\figurefont\mdseries\fontsize{8.8}{10.5}\selectfont,
    text=FigureInk,align=center,text width=28mm] at (0,0)
    {\bfseries 可信机器学习\\[1mm]
     真实使用中的系统\\何时可以被依赖？};

  % 第一层标签。
  \trustlabel[\fontsize{8.2}{9.8}\selectfont]{22.5}{25}{要求与影响}{-67.5}
  \trustlabel[\fontsize{8.2}{9.8}\selectfont]{-112.5}{25}{证据与行为}{-22.5}
  \trustlabel[\fontsize{8.2}{9.8}\selectfont]{-225}{25}{防护与治理}{45}

  % 第二层标签。
  \trustlabel[\fontsize{7.3}{8.8}\selectfont]{67.5}{44.5}{要求\\可操作化}{-22.5}
  \trustlabel[\fontsize{7.3}{8.8}\selectfont]{22.5}{44.5}{未知识别\\与转交}{-67.5}
  \trustlabel[\fontsize{7.3}{8.8}\selectfont]{-22.5}{44.5}{风险、收益\\与负担分配}{67.5}
  \trustlabel[\fontsize{7.3}{8.8}\selectfont]{-67.5}{44.5}{变化、适应\\与恢复}{22.5}
  \trustlabel[\fontsize{7.3}{8.8}\selectfont]{-112.5}{44.5}{解释、审计\\与补救}{-22.5}
  \trustlabel[\fontsize{7.3}{8.8}\selectfont]{-157.5}{44.5}{目标、监督\\与权限}{-67.5}
  \trustlabel[\fontsize{7.3}{8.8}\selectfont]{-202.5}{44.5}{隐私、安全\\与行动保护}{67.5}
  \trustlabel[\fontsize{7.3}{8.8}\selectfont]{-247.5}{44.5}{联合保证\\与持续验证}{22.5}

  % 外层标签：要求与影响。
  \trustlabel[\fontsize{6.5}{7.7}\selectfont]{82.5}{67.5}{危害情景}{-7.5}
  \trustlabel[\fontsize{6.5}{7.7}\selectfont]{67.5}{67.5}{可检验主张}{-22.5}
  \trustlabel[\fontsize{6.5}{7.7}\selectfont]{52.5}{67.5}{处置阈值}{-37.5}
  \trustlabel[\fontsize{6.5}{7.7}\selectfont]{37.5}{67.5}{校准与覆盖}{-52.5}
  \trustlabel[\fontsize{6.5}{7.7}\selectfont]{22.5}{67.5}{拒答与转交}{-67.5}
  \trustlabel[\fontsize{6.5}{7.7}\selectfont]{7.5}{67.5}{主张级核验}{-82.5}

  % 外层标签：证据与行为。
  \trustlabel[\fontsize{6.5}{7.7}\selectfont]{-52.5}{67.5}{复合偏移}{37.5}
  \trustlabel[\fontsize{6.5}{7.7}\selectfont]{-67.5}{67.5}{变化监测}{22.5}
  \trustlabel[\fontsize{6.5}{7.7}\selectfont]{-82.5}{67.5}{适应与恢复}{7.5}
  \trustlabel[\fontsize{6.5}{7.7}\selectfont]{-97.5}{67.5}{归因与反事实}{-7.5}
  \trustlabel[\fontsize{6.5}{7.7}\selectfont]{-112.5}{67.5}{机制解释}{-22.5}
  \trustlabel[\fontsize{6.5}{7.7}\selectfont]{-127.5}{67.5}{补救有效性}{-37.5}

  % 外层标签：防护与治理。
  \trustlabel[\fontsize{6.5}{7.7}\selectfont]{-187.5}{67.5}{数据影响\\与撤回}{82.5}
  \trustlabel[\fontsize{6.5}{7.7}\selectfont]{-202.5}{67.5}{自适应攻击}{67.5}
  \trustlabel[\fontsize{6.5}{7.7}\selectfont]{-217.5}{67.5}{最小权限}{52.5}

  % 统一外轮廓与层级边界。
  \draw[draw=black,line width=.8pt]
    (90:80) arc[start angle=90,end angle=0,radius=80];
  \draw[draw=black,line width=.8pt]
    (-45:80) arc[start angle=-45,end angle=-135,radius=80];
  \draw[draw=black,line width=.8pt]
    (-180:80) arc[start angle=-180,end angle=-225,radius=80];
  \foreach \a in {90,0,-45,-135,-180,-225}
    {\draw[draw=black,line width=.8pt] (\a:55) -- (\a:80);}
  \draw[draw=black,line width=.6pt] (0,0) circle[radius=34];
  \draw[draw=black,line width=.6pt] (0,0) circle[radius=55];
\end{tikzpicture}%
\endgroup
